% !TEX root = ./../main.tex

\chapter{Dynamic Decentralized Aggregation}
\label{chap:decentralized-aggregation}


%%
\section{D3QN Agent}
\label{sec:d3qn-agent}
% Red Q de cada par, replay priorizado y calendario de epsilon (v2v/agents/d3qn_agent.py). Senalar
% donde la implementacion se aparta de la Tabla III del paper (capas ocultas, piso de epsilon).

%%
\section{Centralized Baseline: Federated Averaging}
\label{sec:fedavg}
% Cada 50 slots la BS promedia las redes online de los K agentes y devuelve el promedio
% (v2v/agents/federated.py). Son 2K transferencias por ronda, todas por la BS.

%%
\section{Neighbor-Based Gossip Aggregation}
\label{sec:gossip}
% Cada 50 slots cada agente promedia su red con la de los pares a 300 m o menos, sin la BS
% (v2v/agents/gossip.py), con la variante de a lo sumo S vecinos al azar (--gossip-neighbors).
% Explicar que tiene de dinamica: el grafo de vecinos cambia con las posiciones de cada episodio.

%%
\section{Communication Cost}
\label{sec:communication-cost}
% Transferencias por ronda y carga del nodo mas cargado en cada esquema (comm_models y
% comm_max_load en el CSV). Un modelo pesa 121 588 bytes con K y M por defecto.

\endinput
